\documentclass[11pt]{article}
\usepackage[margin=1in]{geometry}
\usepackage{amsmath,amssymb}
\usepackage{hyperref}
\usepackage{booktabs}

\title{Rapidity frames in $p$Pb and PbPb ultraperipheral collisions:\\
the $\Delta y = 0.465$ shift}
\author{}
\date{\today}

\begin{document}
\maketitle

\section{Summary}

In $p$Pb collisions at the LHC the two beams carry different energies per
nucleon, so the nucleon--nucleon (NN) centre-of-mass frame moves in the
laboratory with rapidity $|\Delta y| = \tfrac12\ln(208/82)\simeq0.465$ towards
the proton. Theory calculations (this code, and Guzey, Innocenti, Sta\'sto
and Strikman, arXiv:2606.05469) are done in the NN centre-of-mass frame,
while detectors measure laboratory rapidities. Comparing the two therefore
requires shifting by $0.465$. In PbPb both beams have the same energy per
nucleon, the centre-of-mass frame is at rest in the laboratory, and no shift
is needed.

\section{Origin of the shift}

The LHC bends both beams with the same magnets on the same orbit, so both
have the same momentum per unit charge (magnetic rigidity). Per nucleon,
this gives
\begin{equation}
  E_{\rm Pb}/{\rm nucleon} = \frac{Z}{A}\,E_p .
\end{equation}
For the 8.16~TeV $p$Pb run: $E_p = 6.5$~TeV and
$E_{\rm Pb}/{\rm nucleon} = 6.5\times82/208 \simeq 2.56$~TeV. For
ultrarelativistic beams,
\begin{equation}
  \sqrt{s_{NN}} = 2\sqrt{E_p\,E_{\rm Pb}} \simeq 2\sqrt{6.5\times2.56}~{\rm TeV}
  \simeq 8.16~{\rm TeV},
\end{equation}
the value used in the code (\texttt{param.ss = 8160} for \texttt{TARGET=pA}).

Take the lead beam along $+z$ and the proton along $-z$. The total light-cone
momenta per nucleon pair are $P^+\propto E_{\rm Pb}$ and $P^-\propto E_p$, so
the rapidity of the NN centre of mass in the laboratory is
\begin{equation}
  Y_{\rm CM} = \frac12\ln\frac{P^+}{P^-} = \frac12\ln\frac{E_{\rm Pb}}{E_p}
  = -\frac12\ln\frac{A}{Z} = -\frac12\ln\frac{208}{82} \simeq -0.465 .
\end{equation}
The centre of mass moves towards the proton, the beam with more momentum per
nucleon.

Rapidity is additive under boosts along the beam axis, so every particle's
rapidity changes by the same constant between the two frames:
\begin{equation}
  y_{\rm lab} = y_{\rm CM} + Y_{\rm CM}.
\end{equation}
The shape of $d\sigma/dy$ is unchanged; the whole distribution is only
translated by $0.465$ units.

\section{Why the calculation lives in the centre-of-mass frame}

The kinematics in the code assume a symmetric NN frame:
\begin{itemize}
  \item The photon energy fraction $z = 2\omega/\sqrt{s_{NN}}$ uses the lead
    Lorentz factor $\gamma = \sqrt{s_{NN}}/(2m_N)$. This is the lead's boost
    in the NN centre-of-mass frame, not in the laboratory.
  \item The $D^0$ rapidity enters through
    $p^+ = m_T e^{y}/\sqrt2$, and the smallest reachable photon energy
    fraction is $z_{\min} = m_T e^{y}/\sqrt{s_{NN}}$
    (\texttt{src/main.cpp}). Both are centre-of-mass quantities.
  \item The momentum fractions probed in the target,
    $x \sim (m_T/\sqrt{s_{NN}})\,e^{-y}$, are likewise centre-of-mass
    expressions.
\end{itemize}
So the $y$ passed to \texttt{D0} and written to the output files is
$y_{\rm CM}$.

\section{Why PbPb needs no shift}

In PbPb both beams have the same $Z/A$, hence the same energy per nucleon
(2.68~TeV each at $\sqrt{s_{NN}}=5.36$~TeV). Then $P^+=P^-$,
$Y_{\rm CM}=0$ and $y_{\rm lab}=y_{\rm CM}$. The same holds for $pp$. The
shift appears only for asymmetric systems ($p$Pb, $p$O, O--Pb, \dots), where
$Z_1/A_1 \neq Z_2/A_2$.

\begin{center}
\begin{tabular}{lccc}
\toprule
System & $E_1$/nucleon & $E_2$/nucleon & $|Y_{\rm CM}|$ \\
\midrule
PbPb, 5.36 TeV & 2.68 TeV & 2.68 TeV & 0 \\
$p$Pb, 8.16 TeV & 6.5 TeV ($p$) & 2.56 TeV (Pb) & $\tfrac12\ln(208/82)\simeq0.465$ \\
\bottomrule
\end{tabular}
\end{center}

\section{Sign conventions}

The size of the shift is fixed; its sign depends on which beam defines
$+y$. There are two conventions to keep apart.

\paragraph{This code and arXiv:2606.05469.} Positive rapidity is the
direction of the photon-emitting lead nucleus; the proton target moves
towards negative rapidity. In the code this follows from
$z_{\min}\propto e^{y}$: larger $y$ requires a more energetic photon, i.e.\
the $D^0$ goes along the photon (lead) direction. The paper states the same
convention explicitly in Sec.~6 and in the captions of its $p$Pb figures
(``the nucleus moves in the positive-rapidity direction and the proton in
the negative-rapidity direction'').

\paragraph{Experiments.} LHC experiments usually take the proton direction
as $+y$ and often quote the two beam configurations ($p$--Pb and Pb--$p$)
separately. Each measurement must be checked for its own convention.

\paragraph{Conversion.} With the proton along $+y$ in the experiment, a
measured laboratory rapidity $y_{\rm lab}^{\rm exp}$ corresponds to
\begin{equation}
  y_{\rm CM}^{\rm exp} = y_{\rm lab}^{\rm exp} - 0.465
  \qquad\text{(proton still along $+y$)},
\end{equation}
and, flipping to the code's orientation (lead along $+y$),
\begin{equation}
  y_{\rm code} = -\left(y_{\rm lab}^{\rm exp} - 0.465\right).
\end{equation}
A data bin $a<y_{\rm lab}^{\rm exp}<b$ therefore corresponds to
$-(b-0.465) < y_{\rm code} < -(a-0.465)$.

\section{Does arXiv:2606.05469 apply the shift?}

No. In the LaTeX source of Guzey, Innocenti, Sta\'sto and Strikman,
``Diffractive open charm photoproduction in ultraperipheral lead-lead and
proton-lead collisions at the LHC'' (arXiv:2606.05469), the $p$Pb section
gives predictions at $\sqrt{s_{NN}}=8.16$~TeV for $-2<y<2$ with the nucleus
along $+y$. The rapidity shift, laboratory versus centre-of-mass frames and
the value $0.465$ are never mentioned. Their $y$ is therefore the NN
centre-of-mass rapidity, the same as ours.

This is consistent: a theory prediction does not need the shift on its own.
Consequences for this project:
\begin{itemize}
  \item Our $p$Pb output (\texttt{run\_pA.sh}) can be compared with their
    $p$Pb figures directly, bin by bin, with no shift, since both use the
    centre-of-mass frame with lead along $+y$.
  \item The shift is needed only when comparing either calculation with
    measured $p$Pb data quoted in laboratory rapidity.
\end{itemize}

\section{Which nucleus emits the photon}

The photon flux scales as $Z^2$, so in $p$Pb UPCs the lead emits essentially
all photons and the proton acts as the target; the proton-emitted
contribution ($Z=1$) is suppressed by $1/82^2$ and neglected. The
distribution is then asymmetric in $y$, which is why the orientation above
matters. In PbPb either nucleus can emit, the measured distribution is
symmetric in $y$, and the physical cross section is the sum of the two
emitter configurations, evaluated at $+y$ and $-y$.

\end{document}
